\chapter{The Double Field Theory Action}\label{ch:dftaction}

The last three chapters assembled a kinematics: doubled coordinates subject to the strong
constraint (\cref{ch:doubled}), a generalized metric $\HH$ that packages the metric $g$ and
the Kalb--Ramond field $b$ into one $O(D,D)$ tensor (\cref{ch:genmetric}), and a gauge algebra
governed by the C-bracket (\cref{ch:courant}). None of this is dynamics yet. The question of
this chapter is: \emph{is there an action, written only with $\HH$, the derivatives
$\partial_M$ and one scalar, whose every index is contracted with an $O(D,D)$-invariant
tensor, and which becomes the familiar low-energy action of the string when nothing depends
on the winding coordinates?} The answer is yes, and the action has four terms. A student
should care for two reasons. First, the construction explains a fact that looks accidental
in \cref{ch:backgrounds}: the relative coefficients $1:4:-\tfrac1{12}$ of $R$,
$(\partial\phi)^2$ and $H^2$ in the NS--NS action\index{NS--NS action} are forced by duality.
Second, the chapter is a clean case study of what the word ``formalized'' can and cannot mean
today: the Lean files attached to this chapter prove statements about integers, and we shall
say precisely how far those are from the action itself.

Throughout, $D$ is the number of spacetime dimensions, indices $i,j,\dots=1,\dots,D$ label
ordinary coordinates, $M,N,\dots=1,\dots,2D$ label doubled ones, and we drop the overall
constant $1/2\kappa_D^2$ in front of every action, as our sources do. The signature is mostly
plus and $\sqrt g$ is shorthand for $\sqrt{|\det g_{ij}|}$.

%=====================================================================================
\section{A dilaton that T-duality leaves alone}\label{sec:dftaction:dilaton}

The low-energy action for the massless NS--NS fields of the closed string is
\begin{equation}\label{eq:dftaction:nsns}
S_{\text{NS}}=\int\dd^Dx\,\sqrt g\,\ee^{-2\phi}\Bigl[R+4(\partial\phi)^2-\tfrac1{12}H_{ijk}H^{ijk}
\Bigr],\qquad H_{ijk}=\partial_ib_{jk}+\partial_jb_{ki}+\partial_kb_{ij},
\end{equation}
where $R$ is the Ricci scalar of $g_{ij}$, $\phi$ is the dilaton\index{dilaton} and $H=\dd b$
is the field strength of the Kalb--Ramond field\index{Kalb--Ramond field}\index{H-flux@$H$-flux}
\source{AMN §3.2, eqs.~(3.17)--(3.18), ll.~868--884; BW, eq.~(7), ll.~343--361}. Its
equations of motion are the conditions for Weyl invariance of the worldsheet theory at one
loop (\cref{ch:backgrounds}).

Look at the measure $\sqrt g\,\ee^{-2\phi}$. Neither factor is invariant under T-duality. Take
one direction to be a circle of radius $R$ and use a coordinate $y\sim y+2\pi\sqrt{\ap}$, so
that $g_{yy}=r^2$ with $r=R/\sqrt{\ap}$ dimensionless. T-duality sends $r\mapsto1/r$, so
$\sqrt g\mapsto\sqrt g/r^2$. But the Buscher rules of \cref{ch:buscher} also shift the dilaton,
$\phi\mapsto\phi-\ln r$, because the string coupling measured by a lower-dimensional observer
must not change. Hence $\ee^{-2\phi}\mapsto r^2\,\ee^{-2\phi}$ and the \emph{product} is
invariant:
\begin{equation}\label{eq:dftaction:buscherinv}
\sqrt g\,\ee^{-2\phi}\;\longmapsto\;\frac{\sqrt g}{r^2}\cdot r^2\ee^{-2\phi}=\sqrt g\,\ee^{-2\phi}.
\end{equation}
This suggests taking the product as the fundamental variable.

\begin{definition}[duality-invariant dilaton]\label{def:dftaction:d}
The field $d$ defined by
\begin{equation}\label{eq:dftaction:d}
\ee^{-2d}=\sqrt g\,\ee^{-2\phi}\qquad\Longleftrightarrow\qquad d=\phi-\tfrac14\ln|\det g|
\end{equation}
is called the $O(D,D)$ dilaton, or shifted dilaton\index{dilaton!O(D,D)@$O(D,D)$ (shifted)}
\index{shifted dilaton}.
\source{AMN §3.3, eq.~(3.24), ll.~956--965}
\end{definition}

Under a global $O(D,D)$ transformation $X'=hX$ the field $d$ is a scalar, $d'(X')=d(X)$, while
$g$ and $b$ transform by the fractional linear rule of \cref{ch:odd}
\source{HHZ §1, eq.~(1.3), ll.~141--149}. The price is that $d$ is not a scalar under
diffeomorphisms: $\ee^{-2d}$ is a \emph{density}\index{density (scalar)}. Under an
infinitesimal diffeomorphism generated by a vector field $\xi^i$ one has
$\delta\phi=\xi^i\partial_i\phi$ and $\delta\sqrt g=\partial_i(\xi^i\sqrt g)$, the latter
because $\delta g_{ij}=\nabla_i\xi_j+\nabla_j\xi_i$ gives
$\delta\sqrt g=\tfrac12\sqrt g\,g^{ij}\delta g_{ij}=\sqrt g\,\nabla_i\xi^i$. Therefore
\begin{equation}\label{eq:dftaction:density}
\delta\bigl(\ee^{-2d}\bigr)=\partial_i\bigl(\xi^i\ee^{-2d}\bigr),\qquad\text{equivalently}
\qquad \delta d=\xi^i\partial_id-\tfrac12\,\partial_i\xi^i .
\end{equation}
The doubled version simply replaces $i$ by $M$: with a gauge parameter
$\xi^M=(\tilde\xi_i,\xi^i)$ that unifies the one-form parameter of $b$-field gauge
transformations with the vector field of diffeomorphisms,
\begin{equation}\label{eq:dftaction:ddensity}
\delta_\xi\bigl(\ee^{-2d}\bigr)=\partial_M\bigl(\xi^M\ee^{-2d}\bigr),\qquad
\delta_\xi d=\xi^M\partial_Md-\tfrac12\,\partial_M\xi^M
\end{equation}
\source{HHZ §1, eq.~(1.8), ll.~216--230; §4.3, eq.~(4.29), ll.~1348--1353; AMN §3.4,
eq.~(3.34), ll.~1100--1112}. A density is exactly what an integration measure must be: if a
Lagrangian is $\ee^{-2d}$ times a gauge scalar $\mathcal S$ (meaning
$\delta_\xi\mathcal S=\xi^M\partial_M\mathcal S$), then
$\delta_\xi(\ee^{-2d}\mathcal S)=\partial_M(\xi^M\ee^{-2d}\mathcal S)$ is a total derivative and
the action is invariant. This one-line argument is the backbone of the proof of gauge invariance
in \cref{sec:dftaction:curvature}.

\begin{physicsbox}[Physical picture: what $d$ measures]
Compactify on a torus of volume $V$. The $D$-dimensional action \eqref{eq:dftaction:nsns}
becomes a lower-dimensional one with prefactor $V\ee^{-2\phi}$: the effective gravitational
coupling of the observer who cannot resolve the torus. T-duality changes the torus and the
ten-dimensional coupling $g_s=\ee^{\phi}$ separately, but it is an exact symmetry of string
perturbation theory, so it cannot change what the lower-dimensional observer measures. The
field $d$ is that invariant coupling, promoted to a local field. String gas cosmology
(\cref{ch:cosmology}) uses the same variable under the name ``shifted dilaton'',
$\varphi=2\phi-\ln V=2d$ \source{BW, eq.~(9), ll.~395--404}.
\end{physicsbox}

%=====================================================================================
\section{Ingredients and index conventions}\label{sec:dftaction:ingredients}

In this chapter we follow the ordering of HHZ and AMN, winding coordinate first:
\begin{equation}\label{eq:dftaction:XM}
X^M=\begin{pmatrix}\tilde x_i\\ x^i\end{pmatrix},\qquad
\partial_M=\begin{pmatrix}\tilde\partial^i\\ \partial_i\end{pmatrix},\qquad
\eta_{MN}=\begin{pmatrix}0&\mathbf 1\\ \mathbf 1&0\end{pmatrix}=\eta^{MN}
\end{equation}
\source{HHZ §1, eqs.~(1.4)--(1.6), ll.~150--188}. Indices are raised and lowered with $\eta$
only, never with $\HH$. With this ordering the generalized metric with \emph{upper} indices is
the matrix that \cref{ch:genmetric} calls $\HH(G,B)$,
\begin{equation}\label{eq:dftaction:Hup}
\HH^{MN}=\begin{pmatrix}g_{ij}-b_{ik}g^{kl}b_{lj}&b_{ik}g^{kj}\\[2pt]
-g^{ik}b_{kj}&g^{ij}\end{pmatrix},\qquad
\HH_{MN}=\eta_{MP}\HH^{PQ}\eta_{QN}=\begin{pmatrix}g^{ij}&-g^{ik}b_{kj}\\[2pt]
b_{ik}g^{kj}&g_{ij}-b_{ik}g^{kl}b_{lj}\end{pmatrix},
\end{equation}
and the second is the inverse of the first, $\HH^{MP}\HH_{PN}=\delta^M{}_N$
\source{HHZ §1, eqs.~(1.9)--(1.11), ll.~256--291}. The only property of the block form that we
need repeatedly is this: \emph{when $\tilde\partial^i=0$, a derivative index $M$ is forced into
the lower block, and the lower-right block of $\HH^{MN}$ is $g^{ij}$.} For instance
\begin{equation}\label{eq:dftaction:box}
\HH^{MN}\partial_M\partial_N f=(g-bg^{-1}b)_{ij}\tilde\partial^i\tilde\partial^jf
+2\,(bg^{-1})_i{}^{j}\,\tilde\partial^i\partial_jf+g^{ij}\partial_i\partial_jf
\;\xrightarrow{\ \tilde\partial=0\ }\;g^{ij}\partial_i\partial_jf .
\end{equation}
In the opposite frame, $\partial_i=0$, the surviving block is $g-bg^{-1}b$; for $b=0$ this is
$g_{ij}$, which is the \emph{inverse} of the dual metric $g'=g^{-1}$ and so plays on
$\tilde x$-space the role that $g^{ij}$ plays on $x$-space. Finally, all fields and gauge parameters obey
the strong constraint\index{strong constraint} of \cref{ch:doubled},
\begin{equation}\label{eq:dftaction:strong}
\eta^{MN}\partial_MA\,\partial_NB=0,\qquad \eta^{MN}\partial_M\partial_NA=0
\qquad\text{for all fields and parameters }A,B,
\end{equation}
so that, up to an $O(D,D)$ rotation, nothing depends on $\tilde x$
\source{HHZ §1, ll.~228--236; AMN §3.3, eqs.~(3.29)--(3.32), ll.~1026--1058}.

\begin{pitfallbox}[Common pitfall: reading the wrong block]
The statement ``$\HH^{MN}\partial_M\partial_N$ reduces to $g^{ij}\partial_i\partial_j$''
depends on the ordering \eqref{eq:dftaction:XM}. With the momentum-first ordering of
\cref{ch:doubled} the same operator is written with the blocks of $\HH$ swapped, and a
careless reader finds $g_{ij}\partial_i\partial_j$, an expression with its indices in
impossible positions. Index position is the safeguard: $\partial_i$ carries a lower index and
can only be contracted with $g^{ij}$.
\end{pitfallbox}

%=====================================================================================
\section{The action}\label{sec:dftaction:action}

\subsection{Which terms can appear?}
We look for a Lagrangian with two derivatives (as in \eqref{eq:dftaction:nsns}), multiplied by
the density $\ee^{-2d}$, built from $\HH$, $d$ and $\partial_M$ with all indices contracted. Two
observations cut the list of candidates down to four.

\emph{The $\Z_2$ of $b\mapsto-b$.} The NS--NS action is invariant under $b\mapsto-b$ (it
contains $b$ only through $H^2$). In the doubled language this reflection acts as
$\tilde\partial^i\mapsto-\tilde\partial^i$ and flips the sign of the off-diagonal blocks of
$\HH$: $\partial\mapsto Z\partial$, $\HH\mapsto Z\HH Z$ with $Z=\diag(-\mathbf 1,\mathbf 1)$.
This $Z$ is \emph{not} an element of $O(D,D)$, because $Z\eta Z=-\eta\neq\eta$. A term in which
every index is contracted directly between $\partial_M$, $\HH^{MN}$ and $\HH_{MN}$ is invariant,
since each contraction produces $ZZ=1$; a term that needs an explicit $\eta$ is not
\source{HHZ §4.1, eqs.~(4.2)--(4.6), ll.~1040--1085}.

\emph{Integration by parts.} Terms with two derivatives on one field, such as
$\HH^{MN}\partial_M\partial_Nd$ and $\partial_M\partial_N\HH^{MN}$, can be traded for products
of first derivatives, because the density $\ee^{-2d}$ produces a factor $-2\partial_Md$ when the
derivative is moved onto it.

What remains is
\begin{equation}\label{eq:dftaction:candidates}
\HH^{MN}\partial_M\HH^{KL}\partial_N\HH_{KL},\qquad
\HH^{MN}\partial_N\HH^{KL}\partial_L\HH_{MK},\qquad
\partial_Md\,\partial_N\HH^{MN},\qquad
\HH^{MN}\partial_Md\,\partial_Nd
\end{equation}
\source{HHZ §4.1, eqs.~(4.7)--(4.9), ll.~1086--1106}. Note that there is no candidate with only
two factors of $\HH$ and no explicit $\eta$: a term like $\partial_M\HH^{KL}\partial^M\HH_{KL}$
needs $\eta^{MN}$ to raise the index, and it vanishes anyway by \eqref{eq:dftaction:strong}.

\subsection{The result}

\begin{theorem}[Hohm--Hull--Zwiebach action]\label{thm:dftaction:action}
\index{double field theory!action}\index{Hohm--Hull--Zwiebach action}
The action
\begin{equation}\label{eq:dftaction:S}
\begin{aligned}
S_{\text{DFT}}=\int\dd^Dx\,\dd^D\tilde x\;\ee^{-2d}\Bigl[\;&\tfrac18\,\HH^{MN}\partial_M\HH^{KL}
\partial_N\HH_{KL}-\tfrac12\,\HH^{MN}\partial_N\HH^{KL}\partial_L\HH_{MK}\\
&-2\,\partial_Md\,\partial_N\HH^{MN}+4\,\HH^{MN}\partial_Md\,\partial_Nd\;\Bigr]
\end{aligned}
\end{equation}
is invariant under global $O(D,D)$ transformations and, when the strong constraint
\eqref{eq:dftaction:strong} holds, under the gauge transformations
$\delta_\xi\HH^{MN}=\widehat{\mathcal L}_\xi\HH^{MN}$ (generalized Lie derivative,
\cref{ch:courant}) and \eqref{eq:dftaction:ddensity}. \TL
\source{HHZ §1, eq.~(1.15), ll.~336--352; §4.1, eqs.~(4.10)--(4.11), ll.~1107--1126}
\end{theorem}

For reference, the generalized Lie derivative\index{generalized Lie derivative} of $\HH$ is
\begin{equation}\label{eq:dftaction:genlie}
\widehat{\mathcal L}_\xi\HH^{MN}=\xi^P\partial_P\HH^{MN}+(\partial^M\xi_P-\partial_P\xi^M)\HH^{PN}
+(\partial^N\xi_P-\partial_P\xi^N)\HH^{MP}
\end{equation}
\source{HHZ §1, eqs.~(1.16)--(1.17), ll.~362--372}: a transport term, and for each index the
difference of a ``covariant'' and a ``contravariant'' rotation, so that $\eta^{MN}$ itself is
invariant.

The global invariance is immediate, and this is the point of the generalized-metric formulation.
Under $X'=hX$ with $h^{\mathsf T}\eta h=\eta$ the fields transform linearly,
$\HH'(X')=h\,\HH(X)\,h^{\mathsf T}$ for the upper-index matrix, $\partial'=(h^{-1})^{\mathsf T}
\partial$ and $d'(X')=d(X)$ \source{HHZ §1, eq.~(1.10), ll.~272--278}. Every index in
\eqref{eq:dftaction:S} is contracted between an upper and a lower position, so all factors of
$h$ cancel in pairs, and the measure is invariant because $|\det h|=1$. If some directions are
compact, the periodicities of $(x,\tilde x)$ must also be preserved and the group is reduced to
the discrete subgroup $\Odd{d}$ of the $d$ toroidal directions \source{HHZ §1, ll.~191--192}.
Compare this with the earlier formulation in terms of $E_{ij}=g_{ij}+b_{ij}$, where the same
symmetry acts by fractional linear transformations and had to be verified term by term
\source{HHZ §1, eqs.~(1.1)--(1.3), ll.~100--149 and ll.~245--252}. Gauge invariance is less
obvious; we return to it in \cref{sec:dftaction:curvature}.

\begin{figure}[ht]
\centering
\begin{tikzpicture}[>=Stealth,node distance=12mm,
 box/.style={draw,rounded corners=2pt,align=center,inner sep=5pt,font=\small}]
\node[box] (dft) {$S_{\text{DFT}}[\HH,d]$ on the doubled space $(x^i,\tilde x_i)$\\
 strong constraint imposed};
\node[box,below left=14mm and -14mm of dft] (sugra) {$\int\dd^Dx\,\sqrt g\,
 \ee^{-2\phi}\bigl[R+4(\partial\phi)^2-\tfrac1{12}H^2\bigr]$\\ fields $g,b,\phi$ on $x$-space};
\node[box,below right=14mm and -14mm of dft] (dual) {the same functional of the\\ dual fields
 $g',b',\phi'$ on $\tilde x$-space};
\draw[->] (dft) -- node[left=2mm,font=\small] {$\tilde\partial^i=0$} (sugra);
\draw[->] (dft) -- node[right=2mm,font=\small] {$\partial_i=0$} (dual);
\draw[<->] (sugra.south east) to[bend right=25] node[below,font=\small]
 {e.g.\ $h=\eta\in O(D,D)$} (dual.south west);
\end{tikzpicture}
\caption{One action, several supergravity frames. Every solution of the strong constraint is an
$O(D,D)$ rotation of the frame $\tilde\partial^i=0$ \source{AMN §3.3, ll.~1050--1056}; because $S_{\text{DFT}}$ is $O(D,D)$
invariant, each frame yields the NS--NS action for the correspondingly rotated fields.}
\label{fig:dftaction:frames}
\end{figure}

%=====================================================================================
\section{Reduction to the NS--NS action}\label{sec:dftaction:reduction}

We now set $\tilde\partial^i=0$ (the supergravity frame\index{supergravity frame}) in
\eqref{eq:dftaction:S} and work out each term. The
$\tilde x$-integral then gives a constant volume factor that we absorb in the normalization.
Write $\mathcal L=\ee^{-2d}(T_1+T_2+T_3+T_4)$ for the four terms in the order displayed.

\paragraph{The dilaton terms.} By the rule stated below \eqref{eq:dftaction:Hup}, the
derivative indices select the block $g^{ij}$:
\begin{equation}\label{eq:dftaction:T34}
T_4=4\,g^{ij}\partial_id\,\partial_jd,\qquad T_3=-2\,\partial_id\,\partial_jg^{ij}.
\end{equation}

\paragraph{The first metric term.} In $T_1$ both derivative indices are lower-block, so
\[
T_1=\tfrac18\,g^{ij}\,\partial_i\HH^{KL}\,\partial_j\HH_{KL}.
\]
The contraction over $K,L$ is a sum
over the four blocks of \eqref{eq:dftaction:Hup}, each block of $\HH^{KL}$ being multiplied
entry by entry with the block of $\HH_{KL}$ in the same position. Set $b=0$ first. The two
diagonal blocks give $\partial_ig_{kl}\partial_jg^{kl}$ and $\partial_ig^{kl}\partial_jg_{kl}$,
which are equal after contraction with the symmetric $g^{ij}$, hence
\begin{equation}\label{eq:dftaction:T1g}
T_1\big|_{b=0}=\tfrac14\,g^{ij}\,\partial_ig^{kl}\,\partial_jg_{kl}.
\end{equation}
Now restore $b$ and collect the terms with two derivatives of $b$ and no undifferentiated
$b$. The upper-right blocks contribute
\[
\partial_i(bg^{-1})_k{}^{l}\;\partial_j(-g^{-1}b)^k{}_l\;\longrightarrow\;-(\partial_ib_{km})\,g^{ml}
g^{kn}\,(\partial_jb_{nl})=-\partial_ib_{km}\,\partial_jb^{km},
\]
where the arrow means that we keep only the part in which both derivatives act on $b$, and
where $b^{km}=g^{kn}g^{ml}b_{nl}$. The lower-left blocks give the same, so
\begin{equation}\label{eq:dftaction:T1b}
T_1\big|_{(\partial b)^2}=\tfrac18\,g^{ij}\cdot2\cdot\bigl(-\partial_ib_{kl}\,\partial_jb^{kl}
\bigr)=-\tfrac14\,\partial_ib_{jk}\,\partial^ib^{jk}.
\end{equation}
The remaining $b$-dependent pieces of $T_1$ contain undifferentiated $b$; they come from
$\partial(bg^{-1}b)$ in the upper-left block and from $\partial g^{-1}$ hidden in $bg^{-1}$, and
they cancel among themselves. This must happen: a constant shift of $b$ is an $O(D,D)$
transformation (the $B$-shift of \cref{ch:odd}) that preserves the frame $\tilde\partial=0$, so
undifferentiated $b$ cannot survive.

\paragraph{The second metric term.} In $T_2=-\tfrac12\HH^{MN}\partial_N\HH^{KL}\partial_L
\HH_{MK}$ the indices $N$ and $L$ are lower-block. For $b=0$ the matrices are block diagonal,
which forces $M$ and then $K$ into the lower block as well, and
\begin{equation}\label{eq:dftaction:T2g}
T_2\big|_{b=0}=-\tfrac12\,g^{ij}\,\partial_jg^{kl}\,\partial_lg_{ik}.
\end{equation}
The terms quadratic in $\partial b$ are found in the same way as for $T_1$ (\cref{ex:dftaction:T2b}):
\begin{equation}\label{eq:dftaction:T2b}
T_2\big|_{(\partial b)^2}=-\tfrac12\,\partial_ib_{jk}\,\partial^jb^{ki}.
\end{equation}

\paragraph{Assembling $H^2$.} Expand the square of the three-form in
\eqref{eq:dftaction:nsns}. There are $3\times3=9$ products; three of them pair a term with
itself and six pair different terms. Relabelling indices,
\begin{equation}\label{eq:dftaction:H2}
H_{ijk}H^{ijk}=3\,\partial_ib_{jk}\,\partial^ib^{jk}+6\,\partial_ib_{jk}\,\partial^jb^{ki},
\qquad\text{so}\qquad
-\tfrac1{12}H^2=-\tfrac14\,\partial_ib_{jk}\partial^ib^{jk}-\tfrac12\,\partial_ib_{jk}
\partial^jb^{ki}.
\end{equation}
The two pieces are exactly \eqref{eq:dftaction:T1b} and \eqref{eq:dftaction:T2b}. The
coefficients $\tfrac18$ and $-\tfrac12$ of the $O(D,D)$ action therefore \emph{produce} the
$\tfrac1{12}$ of supergravity, and they also fix the gauge-invariant combination $H=\dd b$:
neither $T_1$ nor $T_2$ is invariant under $b\mapsto b+\dd\tilde\lambda$ alone.

\begin{proposition}[reduction, first form]\label{prop:dftaction:L0}
For $\tilde\partial^i=0$ the Lagrangian of \eqref{eq:dftaction:S} equals, with no integration by
parts,
\begin{equation}\label{eq:dftaction:L0}
\mathcal L^{(0)}=\ee^{-2d}\Bigl[\tfrac14\,g^{ij}\partial_ig_{kl}\,\partial_jg^{kl}
-\tfrac12\,g^{ij}\partial_jg^{kl}\,\partial_lg_{ik}-2\,\partial_id\,\partial_jg^{ij}
+4\,g^{ij}\partial_id\,\partial_jd-\tfrac1{12}H^2\Bigr].
\end{equation}
\TL \source{HHZ §4.1, eqs.~(4.12)--(4.13), ll.~1127--1170}
\end{proposition}

\noindent We derived every term above except the cancellation of undifferentiated $b$. As an
independent symbolic check we evaluated both sides with sympy in $D=3$ for a non-diagonal
polynomial metric, a generic polynomial $b$ and a non-constant $d$, building $\HH$ from
\eqref{eq:dftaction:Hup}; the difference is identically zero, and the separate identities
\eqref{eq:dftaction:T1b}, \eqref{eq:dftaction:T2b} hold for the full $b$-dependence of $T_1$ and
$T_2$, not only at quadratic order.

\paragraph{From first derivatives to the Ricci scalar.} Equation \eqref{eq:dftaction:L0}
contains no second derivatives, whereas $R$ does. The two differ by a total derivative, just as
the Einstein--Hilbert Lagrangian $\sqrt g\,R$ differs by a total derivative from a Lagrangian
quadratic in Christoffel symbols. The shortest way to see the match uses the scalar
$\mathcal R$ of the next section, so we state the result here and complete the argument there.

\begin{theorem}[reduction to supergravity]\label{thm:dftaction:reduction}
With $\tilde\partial^i=0$, $\HH$ parametrized by \eqref{eq:dftaction:Hup} and $d$ by
\eqref{eq:dftaction:d},
\begin{equation}\label{eq:dftaction:reduction}
\mathcal L_{\text{DFT}}\Big|_{\tilde\partial=0}=\sqrt g\,\ee^{-2\phi}\Bigl[R+4(\partial\phi)^2-
\tfrac1{12}H^2\Bigr]+\partial_i\bigl(\ee^{-2d}\,V^i\bigr),\qquad
V^i=\partial_jg^{ij}+g^{ij}\partial_j\ln|\det g| .
\end{equation}
\TL \source{HHZ §4.1, ll.~1127--1185; AMN §3.6, eq.~(3.71) and ll.~1458--1477}
\end{theorem}

%=====================================================================================
\section{The generalized scalar curvature}\label{sec:dftaction:curvature}

\subsection{Integrating by parts into the form \texorpdfstring{$\ee^{-2d}\mathcal R$}{exp(-2d) R}}
The Einstein--Hilbert Lagrangian is ``density times scalar'', and gauge invariance follows from
that structure alone. The Lagrangian in \eqref{eq:dftaction:S} is not of this form: its four terms
are not separately gauge scalars. We now bring it to that form. Take the two dilaton terms and
use $\partial_M\ee^{-2d}=-2\,\ee^{-2d}\partial_Md$:
\begin{align}
-2\,\ee^{-2d}\partial_Md\,\partial_N\HH^{MN}
 &=\partial_M\bigl(\ee^{-2d}\bigr)\partial_N\HH^{MN}
 =\partial_M\bigl(\ee^{-2d}\partial_N\HH^{MN}\bigr)-\ee^{-2d}\,\partial_M\partial_N\HH^{MN},
 \label{eq:dftaction:ibp1}\\
4\,\ee^{-2d}\HH^{MN}\partial_Md\,\partial_Nd
 &=-4\,\ee^{-2d}\HH^{MN}\partial_Md\,\partial_Nd+8\,\ee^{-2d}\HH^{MN}\partial_Md\,\partial_Nd
 \notag\\
 &=-4\,\ee^{-2d}\HH^{MN}\partial_Md\,\partial_Nd-4\,\partial_M\bigl(\ee^{-2d}\bigr)\HH^{MN}
 \partial_Nd\notag\\
 &=-4\,\ee^{-2d}\HH^{MN}\partial_Md\,\partial_Nd-4\,\partial_M\bigl(\ee^{-2d}\HH^{MN}\partial_Nd
 \bigr)+4\,\ee^{-2d}\,\partial_M\bigl(\HH^{MN}\partial_Nd\bigr).\label{eq:dftaction:ibp2}
\end{align}
Adding the two lines and expanding the last bracket gives
\begin{equation}\label{eq:dftaction:LR}
\mathcal L_{\text{DFT}}=\ee^{-2d}\,\mathcal R+\partial_M\Bigl(\ee^{-2d}\bigl[\partial_N\HH^{MN}
-4\,\HH^{MN}\partial_Nd\bigr]\Bigr),
\end{equation}
with
\begin{equation}\label{eq:dftaction:R}
\begin{aligned}
\mathcal R={}&4\,\HH^{MN}\partial_M\partial_Nd-\partial_M\partial_N\HH^{MN}
-4\,\HH^{MN}\partial_Md\,\partial_Nd+4\,\partial_M\HH^{MN}\partial_Nd\\
&+\tfrac18\,\HH^{MN}\partial_M\HH^{KL}\partial_N\HH_{KL}
-\tfrac12\,\HH^{MN}\partial_M\HH^{KL}\partial_K\HH_{NL}
\end{aligned}
\end{equation}
\index{generalized scalar curvature}\source{HHZ §4.2, eqs.~(4.24)--(4.27), ll.~1272--1338}.
(The last term is the second term of \eqref{eq:dftaction:S} with its dummy indices renamed,
$M\leftrightarrow N$ and then $K\leftrightarrow L$, using the symmetry of $\HH$.)

\begin{theorem}[gauge scalar]\label{thm:dftaction:scalar}
If the strong constraint holds, $\delta_\xi\mathcal R=\xi^M\partial_M\mathcal R$. Consequently
$S_{\text{DFT}}=\int\ee^{-2d}\mathcal R$ (up to a boundary term) is gauge invariant. \TL
\source{HHZ §4.3, eqs.~(4.28)--(4.47), ll.~1340--1528; AMN §3.6, eqs.~(3.66)--(3.67),
ll.~1405--1421}
\end{theorem}

\noindent\emph{Idea of the proof.} Every object $W$ built from $\HH$, $d$ and partial
derivatives transforms as $\delta_\xi W=\widehat{\mathcal L}_\xi W+\Delta_\xi W$, where the first
term is what a generalized tensor with the index structure of $W$ would do and $\Delta_\xi W$
is the failure. Partial derivatives are not covariant, so for example
$\Delta_\xi(\partial_Md)=-\tfrac12\partial_M(\partial\cdot\xi)$: differentiating
\eqref{eq:dftaction:ddensity} gives $\xi^P\partial_P\partial_Md+\partial_M\xi^P\partial_Pd
-\tfrac12\partial_M\partial_P\xi^P$; the first two terms are the generalized Lie derivative of a
covector, because the term $-\partial^P\xi_M\,\partial_Pd$ that would complete it vanishes by
the strong constraint, and the last term is the failure. The map $\Delta_\xi$ obeys the Leibniz rule, so one computes it
on each of the six terms of \eqref{eq:dftaction:R} and discards everything of the form
$\partial^P(\cdot)\,\partial_P(\cdot)$ by the strong constraint. The four terms containing $d$
fail together by
\begin{equation}\label{eq:dftaction:failure}
\Delta_\xi\bigl[4\HH^{MN}\partial_M\partial_Nd-\partial_M\partial_N\HH^{MN}+4\partial_M\HH^{MN}
\partial_Nd-4\HH^{MN}\partial_Md\,\partial_Nd\bigr]=\partial_M\partial_N\xi^P\,\partial_P\HH^{MN},
\end{equation}
all terms with three derivatives of $\xi$ having cancelled inside the bracket, and the two
pure-$\HH$ terms fail by exactly the opposite amount
\source{HHZ §4.3, eqs.~(4.30)--(4.47), ll.~1354--1528}. We do not reproduce the two pages of
algebra. The point to retain is where the strong constraint enters: not in writing
the action, but in its gauge invariance.

\subsection{\texorpdfstring{$\mathcal R$}{R} is the dilaton equation of motion}
Vary $d$ in \eqref{eq:dftaction:S}. The variation hits the prefactor, giving $-2\,\delta d\,
\mathcal L$, and the two explicit dilaton terms:
\[
\delta_d\mathcal L=-2\,\delta d\,\mathcal L+\ee^{-2d}\bigl[-2\,\partial_M\delta d\,\partial_N
\HH^{MN}+8\,\HH^{MN}\partial_Md\,\partial_N\delta d\bigr].
\]
Integrate the bracket by parts, again using $\partial_M\ee^{-2d}=-2\ee^{-2d}\partial_Md$. Writing
$T_\HH$ for the two pure-$\HH$ terms of \eqref{eq:dftaction:S}, the coefficient of
$\ee^{-2d}\delta d$ is
\[
-2\bigl[T_\HH-2\partial d\,\partial\HH+4\HH\partial d\partial d\bigr]
+\bigl[2\partial\partial\HH-4\partial d\,\partial\HH\bigr]
+\bigl[16\HH\partial d\partial d-8\partial\HH\,\partial d-8\HH\partial\partial d\bigr]
\]
in an obvious shorthand, and collecting terms,
\begin{equation}\label{eq:dftaction:deom}
\delta_dS=-2\int\ee^{-2d}\,\mathcal R\,\delta d .
\end{equation}
So the dilaton equation of motion is $\mathcal R=0$ \source{HHZ §4.2, ll.~1272--1294; AMN §3.7,
eq.~(3.79), l.~1522}. In particular the Lagrangian in the form $\ee^{-2d}\mathcal R$ vanishes on every
solution, a property it shares with the string-frame NS--NS action.

\subsection{Reduction of \texorpdfstring{$\mathcal R$}{R} and the end of the reduction proof}
Setting $\tilde\partial=0$ in \eqref{eq:dftaction:R} must give the dilaton equation of motion
of supergravity, and it does:
\begin{equation}\label{eq:dftaction:Rred}
\mathcal R\Big|_{\tilde\partial=0}=R+4\,\nabla^i\nabla_i\phi-4(\partial\phi)^2-\tfrac1{12}H^2,
\end{equation}
to be compared with the supergravity field equation \source{AMN §3.2, eq.~(3.21), ll.~899--907}.
This is an exact identity, without total derivatives. It is the one formula of this chapter
that we did not derive by hand in general: the second derivatives of $g$ hidden in
$-\partial_M\partial_N\HH^{MN}\to-\partial_i\partial_jg^{ij}$ and in
$4g^{ij}\partial_i\partial_jd$ (recall $d=\phi-\tfrac14\ln|\det g|$) assemble into $R$. We
verified it as a symbolic check with sympy in $D=3$ with the same non-diagonal $g$, generic $b$
and $\phi$ as before, and in \cref{sec:dftaction:example} we verify it by hand for a family of
cosmological backgrounds.

Given \eqref{eq:dftaction:Rred}, the theorem follows in two lines. Since
$\nabla_i(\ee^{-2\phi}\nabla^i\phi)=\ee^{-2\phi}[\nabla^2\phi-2(\partial\phi)^2]$,
\begin{equation}\label{eq:dftaction:boxphi}
\sqrt g\,\ee^{-2\phi}\bigl[4\nabla^2\phi-4(\partial\phi)^2\bigr]=4\,\partial_i\bigl(\sqrt g\,
\ee^{-2\phi}g^{ij}\partial_j\phi\bigr)+4\sqrt g\,\ee^{-2\phi}(\partial\phi)^2 ,
\end{equation}
which turns the combination $4\nabla^2\phi-4(\partial\phi)^2$ of the field equation into the
$+4(\partial\phi)^2$ of the action. Inserting \eqref{eq:dftaction:Rred} and
\eqref{eq:dftaction:boxphi} into \eqref{eq:dftaction:LR} gives \eqref{eq:dftaction:reduction}
with $V^i=\partial_jg^{ij}-4g^{ij}\partial_jd+4g^{ij}\partial_j\phi$, and
$-4\partial_jd+4\partial_j\phi=\partial_j\ln|\det g|$ by \eqref{eq:dftaction:d}. The vector $V^i$
does not contain the dilaton: it is the familiar boundary term that relates $\sqrt g\,R$ to the
first-order ``$\Gamma\Gamma$'' Lagrangian of general relativity. We checked
\eqref{eq:dftaction:reduction}, including the sign of each term of $V^i$, symbolically in $D=2$.

\begin{pitfallbox}[Common pitfall: $4(\partial\phi)^2$ versus $-4(\partial\phi)^2$]
The action \eqref{eq:dftaction:nsns} has $+4(\partial\phi)^2$; the scalar
\eqref{eq:dftaction:Rred} has $-4(\partial\phi)^2+4\nabla^2\phi$. Both are correct and they
differ by the total derivative \eqref{eq:dftaction:boxphi}. The ``wrong-sign'' kinetic term of
the dilaton in the string frame is not a ghost: after the Weyl rescaling to the Einstein frame
(\cref{ch:backgrounds}) the kinetic term has the healthy sign. A related trap: $\mathcal R$ is
\emph{not} the Ricci scalar of the $2D$-dimensional metric $\HH_{MN}$. Its name records only that
it is a gauge scalar with two derivatives that reduces to $R+\dots$; HHZ state that they ``do
not have an understanding of these as curvatures'' \source{HHZ §6, ll.~2446--2451}.
\end{pitfallbox}

\subsection{The generalized Ricci tensor}
Varying $\HH$ is more delicate than varying $d$, because $\HH$ is constrained:
$\HH\eta\HH=\eta^{-1}$. Write the unconstrained variation as
$\delta S=\int\ee^{-2d}\,\delta\HH^{MN}\mathcal K_{MN}$. With $S^M{}_N=\HH^{MP}\eta_{PN}$, which
satisfies $S^2=1$, an allowed variation obeys $\delta\HH=-S\,\delta\HH\,S^{\mathsf T}$, and the
projectors $\tfrac12(1\pm S)$ select the part of $\mathcal K$ that such variations can see. The
field equation is
\begin{equation}\label{eq:dftaction:ricci}
\mathcal R_{MN}\equiv\tfrac14(1-S^{\mathsf T})\mathcal K(1+S)+\tfrac14(1+S^{\mathsf T})
\mathcal K(1-S)=0,
\end{equation}
\index{generalized Ricci tensor}which combines the $g$ and $b$ equations of supergravity in one
$O(D,D)$ tensor \source{HHZ §4.4, eqs.~(4.48)--(4.58), ll.~1536--1650}. The counting is
instructive: $\delta\HH$ maps the $+1$ eigenspace of $S$ to the $-1$ eigenspace and back, so it
has $D\times D$ independent components, the $\tfrac12D(D+1)$ of $\delta g$ plus the
$\tfrac12D(D-1)$ of $\delta b$.

%=====================================================================================
\section{Worked example: cosmological backgrounds and scale-factor duality}
\label{sec:dftaction:example}

Let $D=n+1$, let all fields depend on time $t$ only, put $b=0$ and
\begin{equation}\label{eq:dftaction:cosmo}
\dd s^2=-\dd t^2+\sum_{i=1}^na_i(t)^2(\dd x^i)^2,\qquad h_i\equiv\frac{\dot a_i}{a_i},\qquad
\ee^{-2d}=\Bigl(\prod_ia_i\Bigr)\ee^{-2\phi}
\end{equation}
\source{BW, eq.~(8), ll.~373--393}. We write $h_i$ for the Hubble rates to avoid a clash with
the three-form $H$.

\paragraph{The DFT side.} Only $\partial_t$ is non-zero and $\HH^{tt}=g^{tt}=-1$. In
\eqref{eq:dftaction:L0}, the second and third terms contain $\partial_tg^{tt}$ or
$\partial_tg_{tt}$ and vanish. The first term needs
\[
\partial_tg_{kl}\,\partial_tg^{kl}=\sum\nolimits_i(2a_i\dot a_i)(-2\dot a_i/a_i^3)=
-4\sum\nolimits_ih_i^2,
\]
hence $\tfrac14g^{tt}(-4\sum h_i^2)=+\sum h_i^2$; the fourth is $4g^{tt}\dot d^2=-4\dot d^2$. Therefore
\begin{equation}\label{eq:dftaction:Lcosmo}
\mathcal L_{\text{DFT}}=\ee^{-2d}\Bigl[\sum_{i=1}^nh_i^2-4\,\dot d^{\,2}\Bigr].
\end{equation}
This took three lines. Scale-factor duality\index{scale-factor duality}, $a_i\mapsto1/a_i$ for
any subset of directions with $d$ fixed, flips $h_i\mapsto-h_i$ and is manifest
\source{BW, eq.~(10), ll.~405--431}. In DFT language it is the $O(D,D)$ element that exchanges
$x^i\leftrightarrow\tilde x_i$ in those directions; it acts on $\HH=\diag(g,g^{-1})$ by swapping
$a_i^2\leftrightarrow a_i^{-2}$ and leaves \eqref{eq:dftaction:S} invariant term by term.

\paragraph{The supergravity side.} For \eqref{eq:dftaction:cosmo} the Ricci scalar is
$R=2\sum_i\dot h_i+\sum_ih_i^2+(\sum_ih_i)^2$ (standard; for $n=3$ and equal $a_i$ it is the
familiar $6\dot h+12h^2$). From \eqref{eq:dftaction:cosmo}, $\dot\phi=\dot d+\tfrac12\sum_ih_i$,
so $4(\partial\phi)^2=-4\dot\phi^2=-4\dot d^2-4\dot d\sum h_i-(\sum h_i)^2$. The
$(\sum h_i)^2$ terms cancel and
\begin{align*}
\sqrt g\,\ee^{-2\phi}\bigl[R+4(\partial\phi)^2\bigr]
&=\ee^{-2d}\Bigl[2\sum\nolimits_i\dot h_i+\sum\nolimits_ih_i^2-4\dot d^{\,2}
-4\dot d\sum\nolimits_ih_i\Bigr]\\
&=\ee^{-2d}\Bigl[\sum\nolimits_ih_i^2-4\dot d^{\,2}\Bigr]+\frac{\dd}{\dd t}\Bigl(2\,\ee^{-2d}
\sum\nolimits_ih_i\Bigr),
\end{align*}
because $\frac{\dd}{\dd t}(2\ee^{-2d}\sum h_i)=\ee^{-2d}[2\sum\dot h_i-4\dot d\sum h_i]$. This is
\cref{thm:dftaction:reduction} for this family, by hand; one checks that $-\ee^{-2d}V^t$ with
$V^t=g^{tt}\partial_t\ln|\det g|=-2\sum h_i$ is the same boundary term. On the supergravity side
the duality is hidden: $R$ and $\phi$ each change under $a_i\mapsto1/a_i$, and the invariance
appears only after the cancellation of $(\sum h_i)^2$.

\paragraph{Solutions, with numbers.} There are three equations. Varying $\ln a_i$ in
\eqref{eq:dftaction:Lcosmo} gives $\frac{\dd}{\dd t}(\ee^{-2d}h_i)=0$. Varying $d$ gives
$\mathcal R=0$, and evaluating \eqref{eq:dftaction:R} on \eqref{eq:dftaction:cosmo} term by term
(only the first, third, and fifth terms survive) gives
$\mathcal R=-4\ddot d+4\dot d^{\,2}+\sum h_i^2$. The third is the equation of $g_{tt}$, which we
lost by fixing $g_{tt}=-1$: restoring $g_{tt}=-N(t)^2$ one finds
$\mathcal L=\ee^{-2\hat d}N^{-1}[\sum h_i^2-4\dot{\hat d}{}^2]$ with
$\ee^{-2\hat d}=\prod a_i\,\ee^{-2\phi}$ (symbolic check), and varying $N$ at $N=1$ gives the
constraint $\sum h_i^2=4\dot d^{\,2}$. Try the power-law ansatz $a_i=t^{p_i}$, $\ee^{-2d}=t$, for
which $h_i=p_i/t$, $\dot d=-1/2t$ and $\ddot d=1/2t^2$. The first equation holds because
$\ee^{-2d}h_i=p_i$ is constant. The constraint reads $\sum p_i^2/t^2=1/t^2$ and the dilaton
equation reads $(-2+1+\sum p_i^2)/t^2=0$; both hold if and only if
\begin{equation}\label{eq:dftaction:kasner}
\sum_{i=1}^np_i^2=1 .
\end{equation} For $n=3$ isotropic directions, $p=1/\sqrt3\approx0.577$
and the physical dilaton is $\phi=d+\tfrac12\sum p_i\ln t=\tfrac12(\sqrt3-1)\ln t\approx
0.366\ln t$: a decelerated expansion with growing coupling. Its scale-factor dual has
$p=-1/\sqrt3$ and $\phi=-\tfrac12(\sqrt3+1)\ln t\approx-1.366\ln t$: contraction with a
coupling that decreases. The two are one point of the constraint sphere
\eqref{eq:dftaction:kasner} and its antipode; in the DFT variables $(\HH,d)$ they are related by
a constant $O(D,D)$ matrix and share the same $d(t)$.

\begin{figure}[ht]
\centering
\begin{tikzpicture}[>=Stealth,xscale=2.1,yscale=1.25]
\draw[->] (0,0) -- (3.2,0) node[right] {$t$};
\draw[->] (0,0) -- (0,2.6) node[above] {$a(t)$};
\draw[thick,domain=0.02:3,samples=60] plot (\x,{\x^0.577});
\draw[thick,dashed,domain=0.2:3,samples=60] plot (\x,{\x^(-0.577)});
\draw[dotted] (1,0) node[below] {$1$} -- (1,1) -- (0,1) node[left] {$1$};
\node[font=\small] at (2.55,2.05) {$a=t^{1/\sqrt3}$};
\node[font=\small] at (2.55,0.33) {$a=t^{-1/\sqrt3}$};
\end{tikzpicture}
\caption{The isotropic $n=3$ solution of \eqref{eq:dftaction:kasner} and its scale-factor dual.
Both have the same $\ee^{-2d}=t$; they cross at the self-dual value $a=1$, that is
$R=\sqrt{\ap}$ in the units of \cref{sec:dftaction:dilaton}.}\label{fig:dftaction:sfd}
\end{figure}

%=====================================================================================
\section{Beyond the strong constraint: what is and is not known}\label{sec:dftaction:beyond}

It is tempting to read \eqref{eq:dftaction:S} as a theory of fields on a $2D$-dimensional space.
With the strong constraint it is not: the constraint implies that, after an $O(D,D)$ rotation,
no field depends on $\tilde x$, and the action is then the NS--NS action in a particular
notation \source{HHZ §1, ll.~228--236; AMN §3.3, ll.~1050--1058}. What has been gained is
(i) a formulation in which $O(D,D)$ acts linearly and the duality invariance of the
dimensionally reduced theory is visible before reduction, (ii) a geometric unification of
diffeomorphisms and $b$-field gauge transformations, and (iii) a starting point for relaxing
the constraint. The status of (iii), according to the sources of this book, is as follows.

\begin{itemize}[leftmargin=*]
\item \emph{Weak constraint only.}\index{weak constraint} String field theory requires only
$\partial^M\partial_MA=0$ on each field, not on products. A doubled action with this weaker
condition was constructed to cubic order in fluctuations around a flat background; extending
it to higher order needs new structures, including a projector onto the kernel of
$\partial^M\partial_M$, and HHZ call this ``the most important outstanding question''.
\TL \source{HHZ §6, ll.~2454--2467}
\item \emph{Closure constraints.} What gauge invariance really uses is closure of the gauge
algebra and the vanishing of a small number of obstruction tensors. The strong constraint solves
these conditions, but it is not the only solution: the flux formulation of the action keeps the
terms $\Delta_{(SC)}\mathcal R$ that the strong constraint would set to zero, and
generalized Scherk--Schwarz reductions provide configurations that depend on both $x$ and
$\tilde x$ in a controlled way. \TL \source{AMN §3.5, eqs.~(3.44)--(3.46), ll.~1176--1259; §3.6,
ll.~1443--1472; §3.3, ll.~1071--1083; §5.6, ll.~3703--3720}
\item \emph{Link to the worldsheet.} Under the strong constraint the generalized Ricci-flatness
equation \eqref{eq:dftaction:ricci} is the condition for conformal invariance of a doubled
sigma model; whether this persists beyond the strong constraint is left open in AMN. \TL
\source{AMN §7, ll.~5046--5058}
\end{itemize}
Whether a consistent, fully doubled field theory exists to all orders is therefore \TC\ in the
classification of this book. Nothing in later chapters depends on it: the uses we make of
$\HH$ and $d$ for $\KT$ (\cref{ch:k3t2}) and for the dual-scale principle (\cref{ch:dualscale})
are within the strongly constrained theory or concern constant backgrounds.

%=====================================================================================
\section{What the machine has checked}\label{sec:dftaction:lean}

The honest summary is short: \emph{none of the theorems of this chapter is formalized.} There
is no manifold, no derivative, no matrix-valued field and no integral in the two Lean files
attached to it. What they contain is an \emph{arithmetic shadow}\index{arithmetic shadow}: each
physical quantity is replaced by one integer, and the theorems are identities of integer
arithmetic that the corresponding physical statement would imply. We quote them because a
student who opens the repository will find them under names that promise more, and should
learn to read the statement rather than the name. The files are
\path{DoubleFieldTheory/ActionCurvature.lean} (namespace
\lean{DoubleFieldTheory.ActionCurvature}, no imports) and
\path{DoubleFieldTheory/PhysicsDSL.lean} (namespace \lean{DoubleFieldTheory.PhysicsDSL}). We ran
\lean{\#print axioms} on every theorem quoted below: each depends on a subset of
\lean{propext} and \lean{Quot.sound}, and \lean{dsl\_self\_dual\_minimum} on no axiom at all.

\begin{leanbox}[In Lean: the measure and the Lagrangian as integer products]
\begin{leancode}
def DilatonMeasure (sqrt_g e_minus_2phi : Int) : Int :=
  sqrt_g * e_minus_2phi

theorem dilaton_measure_symm (s e : Int) :
    DilatonMeasure s e = DilatonMeasure e s := by ...

def ActionLagrangian (density ricci : Int) : Int :=
  density * ricci

theorem dft_action_nsns_equivalence (density R_geom : Int) :
    ActionLagrangian density R_geom = ActionLagrangian R_geom density := by ...
\end{leancode}
\textbf{Reading.} \lean{DilatonMeasure} is the product of two integers, named after the two
factors of \eqref{eq:dftaction:d}. The first theorem says that this product does not change when
its two arguments are exchanged, i.e.\ $s\cdot e=e\cdot s$ in $\Z$; the second says the same
for \lean{ActionLagrangian}. Both proofs unfold the definition and rewrite with
\lean{Int.mul\_comm}. Neither statement involves T-duality: the invariance
\eqref{eq:dftaction:buscherinv} would be $(s/r^2)\cdot(r^2e)=s\cdot e$, which needs a field, not
$\Z$ (\cref{ex:dftaction:leaninv}). And despite its name the second theorem does not relate two
actions; the equivalence of \cref{thm:dftaction:reduction} is \TL.
\end{leanbox}

\begin{leanbox}[In Lean: the three-term decomposition]
\begin{leancode}
def DFTRicciComponents (R_geom kin_phi H_sq : Int) : Int :=
  R_geom + 4 * kin_phi - H_sq

theorem dft_ricci_expansion (R_geom kin_phi H_sq : Int) :
    DFTRicciComponents R_geom kin_phi H_sq + H_sq = R_geom + 4 * kin_phi := by ...

theorem dft_ricci_physical_reduction (R_geom : Int) :
    DFTRicciComponents R_geom 0 0 = R_geom := by ...

theorem connection_trace_conservation (trace_H_initial delta_trace : Int)
    (hd : delta_trace = 0) :
    trace_H_initial + delta_trace = trace_H_initial := by ...
\end{leancode}
\textbf{Reading.} \lean{DFTRicciComponents} is the integer $R+4k-h$. To match the bracket of
\eqref{eq:dftaction:nsns} one must read $k$ as $(\partial\phi)^2$ and $h$ as
$\tfrac1{12}H^2$, not as $H^2$: the coefficient $\tfrac1{12}$, which
\cref{sec:dftaction:reduction} showed to be the non-trivial output of the reduction, is absent
from the file. The two theorems say $(R+4k-h)+h=R+4k$ and $R+4\cdot0-0=R$; both are closed by
the linear-arithmetic tactic \lean{omega} after unfolding. For example
$\lean{DFTRicciComponents}\ 5\ 3\ 2=15$. The third theorem says $a+\delta=a$ when $\delta=0$; its
docstring relates it to compatibility of a connection with $\HH$, but no connection occurs in
the statement. What would be a faithful version of the second theorem? It is the statement that
\eqref{eq:dftaction:Rred} equals $R$ when $\phi$ is constant and $H=0$, a theorem about fields.
\end{leanbox}

\begin{leanbox}[In Lean: the trace of the Einstein tensor, as arithmetic]
\begin{leancode}
def ContractedEinstein (D R : Int) : Int :=
  (D - 2) * R

theorem contracted_einstein_dim2 (R : Int) :
    ContractedEinstein 2 R = 0 := by ...
\end{leancode}
\textbf{Reading.} $(2-2)\cdot R=0$ for every integer $R$. The physics behind the name: the
Einstein tensor $G_{\mu\nu}=R_{\mu\nu}-\tfrac12g_{\mu\nu}R$ has trace
$g^{\mu\nu}G_{\mu\nu}=R-\tfrac D2R=-\tfrac12(D-2)R$, which vanishes identically on a
two-dimensional worldsheet (\cref{ch:classical}). The Lean definition drops the factor
$-\tfrac12$, which is harmless for the vanishing at $D=2$ and wrong for any other use; e.g.\
$\lean{ContractedEinstein}\ 4\ 7=14$ whereas the trace is $-7$.
\end{leanbox}

\begin{leanbox}[In Lean: the notation layer and the one-component section condition]
\begin{leancode}
structure FieldDeriv where
  dx : Int
  dtx : Int

def SectionContract (Phi Psi : FieldDeriv) : Int :=
  Phi.dx * Psi.dtx + Phi.dtx * Psi.dx

notation:65 "∂_M " Φ " ∂^M " Ψ => SectionContract Φ Ψ

theorem dsl_strong_section_condition (Φ Ψ : FieldDeriv)
    (hΦ : Φ.dtx = 0) (hΨ : Ψ.dtx = 0) :
    (∂_M Φ ∂^M Ψ) = 0 := by ...

def Reff (R : Nat) : Nat :=
  R * R + 1

theorem dsl_effective_radius_strictly_super_planckian (R : Nat) (h : R ≥ 1) :
    R_eff(R) ≥ 2 := by ...
\end{leancode}
\textbf{Reading.} The first two declarations live in
\path{DoubleFieldTheory/CourantAlgebroid.lean}; \path{PhysicsDSL.lean} adds \emph{notation} so
that Lean source resembles the formulas of a physics paper, and restates earlier theorems in that
notation (\lean{dsl\_courant\_pairing\_symm}, \lean{dsl\_cbracket\_antisymm},
\lean{dsl\_cbracket\_self\_vanishes}, \lean{dsl\_dorfman\_cbracket\_diff}; see
\cref{ch:courant}). A \lean{FieldDeriv} is a pair of integers standing for the values of
$\partial_xA$ and $\tilde\partial^xA$ of one field at one point for $D=1$, and
\lean{SectionContract} is $\eta^{MN}\partial_MA\,\partial_NB$ for such pairs. The theorem says:
if the second entries of both pairs vanish, the contraction vanishes. This is one direction only
(supergravity frame $\Rightarrow$ strong constraint), for $D=1$, at one point. The all-$d$
statement on the charge lattice, including $O(d,d;\Z)$ covariance, is
\lean{IsSection} and \lean{isSection\_image} of \cref{ch:doubled}. Finally \lean{Reff R} is
$R^2+1$, the numerator of $R+1/R$ in units $\ap=1$, and the last theorem says $R^2+1\ge2$ for
natural numbers $R\ge1$. This is weaker than the dual-scale inequality $R+1/R\ge2$, which would
read $R^2+1\ge2R$; the faithful real-number statement is
\lean{circle\_effective\_scale\_ge\_two} in \path{DualScaleStream2/DualScale/TraceBound.lean}
(\cref{ch:dualscale}). The docstring's conclusion about singularities is an interpretation,
\TC\ at best, and not part of what is proved.
\end{leanbox}

\begin{pitfallbox}[Common pitfall: trusting docstrings]
The Lean kernel checks statements and proofs. It does not read comments. The docstrings of
\path{ActionCurvature.lean} contain correct physics (they quote \eqref{eq:dftaction:S} and
\eqref{eq:dftaction:reduction}), equation numbers of papers that we have not checked against the
papers and that do not match the numbering of HHZ used in this chapter, and claims of
verification placed next to one-line integer identities. When you cite a formal
result, cite the statement after the colon, never the prose above it.
\end{pitfallbox}

\begin{table}[ht]
\centering
\caption{Epistemic status of the main claims of this chapter.}\label{tab:dftaction:tiers}
\small
\begin{tabular}{@{}p{0.50\textwidth}cp{0.37\textwidth}@{}}\toprule
Claim & Tier & Basis\\\midrule
$s\cdot e=e\cdot s$, $(R+4k-h)+h=R+4k$, $(2-2)R=0$ in $\Z$ & \TA &
 \lean{dilaton\_measure\_symm}, \lean{dft\_ricci\_expansion}, \lean{contracted\_einstein\_dim2}\\
$\tilde\partial A=\tilde\partial B=0\Rightarrow\eta^{MN}\partial_MA\partial_NB=0$ (integers, $D=1$)
 & \TA & \lean{dsl\_strong\_section\_condition}\\
$\ee^{-2d}=\sqrt g\,\ee^{-2\phi}$ is an $O(D,D)$ scalar and a gauge density & \TL & AMN §3.3--3.4;
 HHZ §1, §4.3\\
Action \eqref{eq:dftaction:S}: $O(D,D)$ and gauge invariance & \TL & HHZ §1, §4\\
Reduction \eqref{eq:dftaction:L0}, \eqref{eq:dftaction:reduction}, \eqref{eq:dftaction:Rred} & \TL
 & HHZ §4.1--4.2, AMN §3.6; symbolic checks in $D=2,3$\\
Pieces \eqref{eq:dftaction:T34}--\eqref{eq:dftaction:H2}; cosmological example & --- & derived on
 the page; not formalized\\
Identification of any Lean object above with $\mathcal R$, $d$ or $S_{\text{DFT}}$ & --- & a
 reading; never \TA\\
A consistent DFT with only the weak constraint, to all orders & \TC & open (HHZ §6)\\
\bottomrule
\end{tabular}
\end{table}

\subsection{What a faithful formalization would need}\label{sec:dftaction:faithful}
It is useful to separate three levels, in increasing order of difficulty.

\emph{Level 1: algebra at a point.} Everything in \cref{sec:dftaction:reduction} is polynomial
algebra once the first derivatives are treated as independent variables. Fix a commutative ring,
an invertible symmetric matrix $g$, and for each $i$ a symmetric matrix $\gamma_i$ (standing for
$\partial_ig$) and an antisymmetric matrix $\beta_i$ (standing for $\partial_ib$). Add scalars $\delta_i$
for $\partial_id$, define $\partial_ig^{-1}:=-g^{-1}\gamma_ig^{-1}$ and the derivative of each block of
\eqref{eq:dftaction:Hup} by the Leibniz rule. Then \cref{prop:dftaction:L0} becomes an identity
between matrix polynomials, of the same kind as the theorems about \lean{genMetric} proved in
\path{DualScaleStream2/DFT/GeneralizedMetric.lean} (\cref{ch:genmetric}), for instance
\lean{tduality\_inverts\_metric}. It needs only Mathlib's \lean{Matrix} library, block matrices
and traces. This would be a genuine \TA\ statement about the four coefficients
$\tfrac18,-\tfrac12,-2,4$, and it is the natural next step (\cref{ex:dftaction:leanjet} is its
$D=1$ case).

\emph{Level 2: fields on $\R^{2D}$.} To state \cref{thm:dftaction:scalar} one needs smooth maps
$\R^{2D}\to$ matrices, partial derivatives and their commutation, the strong constraint as a
hypothesis on a family of functions, and the generalized Lie derivative. Mathlib's calculus
(\lean{fderiv}, iterated derivatives, symmetry of second derivatives) suffices in principle; the
obstacle is the volume of index bookkeeping, since Mathlib has no tensor-index notation and each
contraction must be written as a finite sum. Integration by parts as in
\eqref{eq:dftaction:LR} would be stated as an identity of integrands (``equal up to a
divergence''), avoiding measure theory.

\emph{Level 3: the geometric statement.} \Cref{thm:dftaction:reduction} mentions the Ricci
scalar of a pseudo-Riemannian metric. A formal version on a manifold needs metrics of
indefinite signature, the Levi-Civita connection, curvature, densities and differential forms
with $\dd b$. To our knowledge these were not all available in Mathlib when the libraries of
this book were compiled; we have not audited Mathlib for this chapter, so treat this sentence
as a pointer, not a fact. On $\R^D$ with global coordinates one can define $R$ by the coordinate
formula and bypass manifolds, which reduces Level 3 to Level 2.

\begin{summarybox}
\begin{itemize}[leftmargin=*]
\item The combination $\ee^{-2d}=\sqrt g\,\ee^{-2\phi}$ is invariant under T-duality; $d$ is an
 $O(D,D)$ scalar and $\ee^{-2d}$ a density under generalized diffeomorphisms.
\item The DFT action \eqref{eq:dftaction:S} has four terms built from $\HH$, $d$ and
 $\partial_M$ without explicit $\eta$; $O(D,D)$ invariance is manifest, gauge invariance requires
 the strong constraint.
\item With $\tilde\partial=0$ it reduces to \eqref{eq:dftaction:L0} exactly and to the NS--NS
 action up to a dilaton-independent boundary term; the coefficients $\tfrac18$ and $-\tfrac12$
 generate the $-\tfrac1{12}H^2$.
\item Up to a total derivative $\mathcal L=\ee^{-2d}\mathcal R$ with $\mathcal R$ a gauge scalar,
 equal to the dilaton equation of motion; the $\HH$ equation defines a generalized Ricci tensor.
 Neither is known to be a curvature in the Riemannian sense.
\item For cosmological backgrounds the action is
 $\ee^{-2d}[\sum h_i^2-4\dot d^2]$, and scale-factor duality is manifest.
\item With the strong constraint DFT is a reformulation of supergravity; a fully doubled theory
 beyond cubic order is open.
\item Lean: only integer-arithmetic shadows are kernel-checked. A matrix-polynomial version of
 the reduction is within reach of Mathlib and would be the first faithful \TA\ result here.
\end{itemize}
\end{summarybox}

\section*{Exercises}
\addcontentsline{toc}{section}{Exercises}

\begin{exercise}[$\star$]
Show that the two forms of \eqref{eq:dftaction:ddensity} are equivalent. Then set
$\tilde\partial=0$ and $\xi^M=(\tilde\xi_i,0)$, a pure $b$-field gauge transformation, and
conclude that $d$, hence $\phi$, is inert, as it must be.
\end{exercise}

\begin{exercise}[$\star$]
Verify that $Z=\diag(-\mathbf 1,\mathbf 1)$ satisfies $Z\eta Z=-\eta$, and that
$\HH^{MN}\mapsto(Z\HH Z)^{MN}$ is the map $b\mapsto-b$ in \eqref{eq:dftaction:Hup}. Which of the
following are even under this $\Z_2$: $\HH^{MN}\partial_Md\,\partial_Nd$;
$\eta^{MN}\partial_Md\,\partial_Nd$; $\partial_M\HH^{MN}\partial_Nd$?
\end{exercise}

\begin{exercise}[$\star\star$]\label{ex:dftaction:T2b}
Derive \eqref{eq:dftaction:T2b}. Hint: at quadratic order in $\partial b$ and zeroth order in
undifferentiated $b$ you may take the undifferentiated factor $\HH^{MN}$ block diagonal, so
that $M$ and $N$ are both lower-block; both $\partial_N\HH^{KL}$ and $\partial_L\HH_{MK}$ must then
be taken from off-diagonal blocks, with $K$ in the upper block.
Then prove \eqref{eq:dftaction:H2} by expanding the nine products.
\end{exercise}

\begin{exercise}[$\star\star$]
Fill in the algebra leading to \eqref{eq:dftaction:deom}: list the seven terms, and check that
the coefficient of each of the six structures of \eqref{eq:dftaction:R} is $-2$ times the one
displayed there.
\end{exercise}

\begin{exercise}[$\star\star$] \emph{Constant $H$-flux on a flat three-torus.}
Let $D=3$, $g_{ij}=\delta_{ij}$, $\phi=0$ and $b_{12}=-b_{21}=h\,x^3$ with $h$ constant. Compute
$H_{123}$, $H_{ijk}H^{ijk}$ and $\mathcal R$ from \eqref{eq:dftaction:Rred}. (Answer:
$\mathcal R=-h^2/2$.) Is this background a solution? Compute $\HH^{MN}$ and observe that it depends
on $x^3$ only through a $B$-shift matrix; relate this to the monodromy $x^3\to x^3+2\pi$.
\end{exercise}

\begin{exercise}[$\star\star$] \emph{Linear dilaton.}
Let $g_{ij}=\eta_{ij}$ be flat Minkowski space, $b=0$ and $\phi=v_ix^i$ with constant $v$. Show
that $d=\phi$, that $\mathcal R=-4\,v\cdot v$ directly from \eqref{eq:dftaction:R}, and that the
$\HH$ equation is satisfied. Conclude that the background solves the DFT equations if and only if
$v$ is null. What replaces this condition in the noncritical string, where
\eqref{eq:dftaction:nsns} contains an extra constant $c$ \source{BW, eq.~(7), ll.~343--364}?
\end{exercise}

\begin{exercise}[$\star\star$]
Show that the trace of the Einstein tensor in $D$ dimensions is $-\tfrac12(D-2)R$. State
precisely which integer identity \lean{contracted\_einstein\_dim2} proves, and write a corrected
definition over $\Q$ whose value is the trace for every $D$. Does your corrected version still
make the $D=2$ theorem provable by \lean{simp}?
\end{exercise}

\begin{exercise}[$\star$, Lean]\label{ex:dftaction:leaninv}
State and prove in Lean 4 with Mathlib the invariance \eqref{eq:dftaction:buscherinv}:
for real $r\neq0$ and any real $w$, $r^{-1}\cdot(r^2w)=r\,w$. (Here $r=\sqrt g$ for the circle
and $w=\ee^{-2\phi}$; the left side is $\sqrt{g'}\,\ee^{-2\phi'}$.) One tactic suffices:
\lean{field\_simp}. Explain why the statement cannot even be formulated for the integer-valued
\lean{DilatonMeasure}.
\end{exercise}

\begin{exercise}[$\star\star\star$, Lean]\label{ex:dftaction:leanjet}
\emph{Level 1 in $D=1$.} Let $g\neq0$ and $\gamma$ be real numbers, $\gamma$ standing for
$\partial_xg$, so that $\partial_xg^{-1}=-\gamma/g^2$. With $\HH^{MN}=\diag(g,g^{-1})$ show on
paper that the first term of \eqref{eq:dftaction:S} is
$\tfrac18g^{-1}\cdot2\gamma(-\gamma/g^2)=-\tfrac14g^{-1}(\gamma/g)^2$ and that the second
vanishes. Prove the displayed identity in Lean (\lean{field\_simp} followed by \lean{ring}). Then
formulate the general statement for $d\times d$ real matrices using \lean{Matrix.fromBlocks} and
\lean{Matrix.trace}, following the style of \lean{genMetric}; proving it is a small research
project.
\end{exercise}

\section*{Further reading}
\addcontentsline{toc}{section}{Further reading}
\begin{itemize}[leftmargin=*]
\item The action in terms of $\HH$ and $d$, its relation to the earlier $E_{ij}$ formulation and
 the overview of results: \source{HHZ §1, ll.~66--400}. Construction of the action, the $\Z_2$
 argument and the reduction: \source{HHZ §4.1, ll.~1007--1266}. Generalized scalar curvature and
 proof of gauge invariance: \source{HHZ §4.2--4.3, ll.~1268--1528}. Generalized Ricci tensor:
 \source{HHZ §4.4, ll.~1530--1650}. Open problems: \source{HHZ §6, ll.~2411--2470}.
\item The NS--NS action and its equations of motion, the fields $\HH$ and $d$, the strong
 constraint: \source{AMN §3.2--3.3, ll.~850--1084}. Generalized Lie derivative and the density
 $\ee^{-2d}$: \source{AMN §3.4, ll.~1085--1125}. Closure constraints: \source{AMN §3.5,
 ll.~1173--1259}. Frame (flux) formulation of the action and its reduction to the
 generalized-metric form: \source{AMN §3.6, ll.~1260--1477}.
\item Scale-factor duality and the shifted dilaton in cosmology: \source{BW, ll.~340--445}.
\item In this book: the NS--NS action from worldsheet Weyl invariance in \cref{ch:backgrounds};
 the dilaton shift in \cref{ch:buscher}; the strong constraint in \cref{ch:doubled}; $\HH$ and
 its Lean formalization in \cref{ch:genmetric}; generalized Lie derivatives and the C-bracket in
 \cref{ch:courant}; string gas cosmology in \cref{ch:cosmology}; a catalogue of which Lean
 libraries are faithful and which are arithmetic shadows in \cref{ch:architecture}.
\end{itemize}
